\section{Introduction}
\label{sec:introduction}

Language models are increasingly used to generate hypotheses, execute experiments, and evaluate research outputs. End-to-end systems such as the AI Scientist connect these activities into a single pipeline \citep{lu2026towards}. The performance of such a pipeline depends on its research harness: the editable rules, operators, memory, evaluation logic, and stopping conditions that shape what the model does next. A harness can improve without any change to the underlying language model, but demonstrating that improvement requires controls for the proposal space, execution budget, and information available to each condition.

Random legal editing offers a particularly useful reference. If a model-guided process cannot exceed random proposals under the same atomic schema and editable surface, the measured gain may reflect additional search rather than informed adaptation. A fixed initial harness provides a complementary reference by measuring the value of continuation itself. These comparisons become informative only when their champions are evaluated on common support and when failed proposals remain in the end-to-end accounting. Production-scale AutoResearch reports infrastructure fragility, agent memory decay, search-direction stagnation, iteration-cost asymmetry, and metric fixation \citep{chandran2026production}.

This paper asks whether model-guided adaptation improves a complete factor-research process under a fixed model, proposal interface, decision count, and candidate evaluation budget, and whether any effect can be assigned to the information supplied to the proposer. Each model call is asked to return one atomic proposal containing two to five edits. Online execution uses probe information only, while an active branch carries the continuation state and a frozen champion enters one final comparison. Three information conditions vary scalar feedback, cumulative probe memory, and compact diagnostics. Random legal edits and the fixed harness provide matched references across six fresh tasks, two development seeds, and two independent model-call replicates.

The comparison separates quantities that a final score alone would conflate: model calls, expression-candidate positions, and final common support. It also separates end-to-end adaptation from the marginal value of cumulative memory and diagnostics. Failure and call-level variability remain in the process comparison. Novelty and duplicate rate describe select search output, while date-block concentration provides an auxiliary diagnostic for the component contrasts.

The memory-equipped process exceeded both references on the development panel, with median paired gains of 0.003534 over random legal edits and 0.010977 over the fixed harness. The component ablations did not establish stable independent increments from cumulative memory or diagnostics. Median paired differences on the select evaluation showed lower novelty and a higher duplicate rate for memory relative to scalar feedback and random legal edits. These findings support model-guided harness adaptation within the registered setting and show why process evaluation cannot be reduced to a single endpoint.
